% =============================================================================
% Construyendo un cerebro digital para la investigación en libre competencia:
% una guía práctica con Graphify y RAG vectorial para el caso chileno
%
% Documento de investigación / guía práctica — formato sobrio en blanco y negro.
%
% Compilación (desde la raíz del repositorio):
%   latexmk -pdf -output-directory=paper/build paper/main.tex
%
% Limpiar auxiliares de build/ tras compilar (opcional):
%   rm -f paper/build/*.aux paper/build/*.bbl paper/build/*.blg \
%          paper/build/*.fdb_latexmk paper/build/*.fls paper/build/*.log \
%          paper/build/*.out paper/build/*.run.xml paper/build/*.toc \
%          paper/build/*.bib paper/build/main.synctex.gz
% =============================================================================

\documentclass[12pt,a4paper]{article}

% ----- Codificación y lengua -----
\usepackage[utf8]{inputenc}
\usepackage[T1]{fontenc}
\usepackage[spanish,es-nodecimaldot]{babel}

% ----- Tipografía -----
\usepackage{lmodern}
\usepackage{microtype}
% Holgura de emergencia para evitar desbordes horizontales en párrafos con
% tokens largos (rutas, nombres de archivo) sin recurrir a \sloppy global.
\emergencystretch=3em

% ----- Geometría -----
\usepackage[top=3cm,bottom=3cm,left=3.5cm,right=3.5cm]{geometry}

% ----- Matemáticas -----
\usepackage{amsmath}
\usepackage{textcomp}

% ----- Gráficos y colores -----
\usepackage{graphicx}
\usepackage[dvipsnames,table]{xcolor}
% Paleta sobria en blanco y negro. Se conservan los nombres de color para no
% alterar los usos a lo largo del documento; todos apuntan a negro/gris.
\definecolor{darkblue}{RGB}{0,0,0}       % antes azul; ahora negro (títulos, frames)
\definecolor{lightgray}{RGB}{244,244,244} % fondo tenue de recuadros
\definecolor{tabhead}{RGB}{229,229,229}   % cabecera de tabla: gris claro, texto negro
\definecolor{accent}{RGB}{0,0,0}          % antes rojo; ahora negro

% ----- Tablas -----
\usepackage{booktabs}
\usepackage{longtable}
\usepackage{array}
\usepackage{colortbl}
\usepackage{tabularx}

% ----- Marcador sobrio de contenido pendiente (sin color) -----
% Marca discreta en gris para lo que el autor debe completar (p.ej. figuras).
\newcommand{\pend}[1]{\textcolor{black!55}{\textbf{[pendiente: #1]}}}

% ----- Recuadros -----
\usepackage{tcolorbox}
\tcbuselibrary{skins,breakable}
\newtcolorbox{infobox}[1][]{
  colback=lightgray, colframe=darkblue!60,
  fonttitle=\bfseries\small, title=#1,
  breakable, sharp corners
}
\newtcolorbox{keybox}[1][]{
  enhanced, colback=darkblue!4, colframe=darkblue,
  fonttitle=\bfseries, title=#1, coltitle=white,
  attach boxed title to top left={xshift=6mm,yshift=-3mm},
  boxed title style={colback=darkblue, sharp corners},
  breakable, sharp corners
}
% codebox: contenido literal (markdown, comandos). Verbatim para no interpretar
% #, _, %, $, & que aparecen en los anexos (CLAUDE.md, prompts de consulta).
\usepackage{fvextra}  % extiende fancyvrb: breaklines/breakanywhere en Verbatim
\newenvironment{codebox}%
  {\VerbatimEnvironment
   \begin{tcolorbox}[colback=black!3,colframe=black!30,breakable,sharp corners]%
   \begin{Verbatim}[fontsize=\small,breaklines=true,breakanywhere=true]}%
  {\end{Verbatim}\end{tcolorbox}}

% ----- Listas y enumeraciones -----
\usepackage{enumitem}
\setlist[itemize]{noitemsep, topsep=4pt}
\setlist[enumerate]{noitemsep, topsep=4pt}

% ----- Encabezados y secciones -----
\usepackage{titlesec}
\titleformat{\section}{\large\bfseries\color{darkblue}}{}{0em}{}[\titlerule]
\titleformat{\subsection}{\normalsize\bfseries\color{darkblue!80}}{}{0em}{}
\titleformat{\subsubsection}{\normalsize\itshape}{}{0em}{}

% ----- Hyperref (siempre al final, antes de biblatex) -----
\usepackage[colorlinks=true,linkcolor=darkblue,citecolor=darkblue,
            urlcolor=darkblue,pdftitle={Cerebro Digital Libre Competencia Chile}]{hyperref}

% ----- Notas al pie -----
\usepackage{footmisc}

% ----- Comillas tipográficas (requerido por biblatex con babel) -----
\usepackage{csquotes}

% ----- Referencias -----
\usepackage[backend=bibtex,style=footnote-dw,autocite=footnote]{biblatex}
\addbibresource{references.bib}

% ----- Metadatos -----
\title{%
  \textbf{Construyendo un cerebro digital\\
  para la investigación en libre competencia}\\[0.4em]
  \large\textit{Una guía práctica con Graphify y RAG vectorial para el caso chileno}
}
\author{%
  {Sebastián Chacón Salinas} \\
  %\small Facultad de Economía y Negocios, Universidad de Chile\\[0.2em]
  %\small \texttt{schacon@fen.uchile.cl}
}
\date{Junio 2026\\[0.5em]
  \small \textcolor{blue}{Documento de trabajo. Comentarios bienvenidos}}

% =============================================================================
\begin{document}
% =============================================================================

\maketitle
\thispagestyle{empty}

\begin{abstract}
\noindent
Un \emph{cerebro digital} es una infraestructura personal de conocimiento
construida sobre documentos que el investigador selecciona. Mapea las
conexiones entre fuentes, expone patrones e inconsistencias, y produce respuestas
trazables al material subyacente.

Este artículo documenta la adaptación de la metodología de Schrepel\autocite{schrepel2026brain} al derecho de la libre
competencia chileno. Sobre un \textit{corpus} de 520 documentos (sentencias del TDLC y
la Corte Suprema, resoluciones, acuerdos extrajudiciales, instrucciones de
carácter general, expedientes de recomendación normativa y normativa de base) se
construye un \textbf{grafo de conocimiento} que
mapea la estructura del campo, identifica qué conceptos son centrales, cómo evoluciona la
jurisprudencia, dónde hay contradicciones entre niveles de autoridad y qué temas
carecen de desarrollo jurisprudencial propio. 

Adicionalmente, se documenta la generación de un \textbf{sistema RAG vectorial} sobre el mismo \textit{corpus}, que indexa el texto
completo y genera respuestas con cita literal al considerando exacto.

El \textit{pipeline} es replicable, de código abierto y neutro respecto del tamaño de la base de datos. El repositorio está disponible públicamente con el código completo y esta guía para su implementación.
\end{abstract}

\newpage
\tableofcontents
\newpage

% =============================================================================
\section{Introducción}
% =============================================================================

Las autoridades de competencia acumulan jurisprudencia durante décadas. El cuerpo resultante es rico en sustancia pero opaco en estructura. Las conexiones entre documentos son difíciles de captar y no es trivial
rastrear la evolución de las posiciones de las autoridades a lo largo del tiempo.

Este artículo describe un \textit{pipeline} para construir lo que Schrepel llama un
\emph{cerebro digital} orientado al sistema de libre competencia chileno. El punto de
partida es una carpeta de PDFs; el resultado son dos instrumentos de consulta.
El primero, un \textbf{grafo de conocimiento}, hace visible la estructura del
\textit{corpus}, identifica cuáles son los conceptos jurídicos más citados, qué documentos
están aislados y qué conexiones esperadas no existen. El segundo, un \textbf{sistema
RAG vectorial}, permite recuperar el texto exacto de cualquier considerando
relevante a una pregunta planteada en lenguaje natural.

\begin{keybox}[Dos instrumentos para dos preguntas distintas]
\textbf{Grafo (Graphify)} - responde preguntas sobre la \emph{estructura} del
campo: qué es lo más citado, cuáles son los silencios, dónde hay tensión entre
autoridades. No devuelve el texto exacto de los documentos.

\medskip
\textbf{RAG vectorial (ragnar + DuckDB)} - responde preguntas sobre el
\emph{contenido}: qué dijo exactamente la CS en el considerando 19 de tal fallo,
qué fallos mencionan un criterio específico. Produce fragmentos con cita literal.
\end{keybox}

\subsection{Herramientas y equipamiento necesario}

Todas las herramientas del \textit{pipeline} son de código abierto y gratuitas,
excepto la suscripción a Claude Code.

\begin{table}[htbp]
\centering
\caption{Herramientas necesarias para replicar el \textit{pipeline} completo}
\label{tab:herramientas}
\small
\begin{tabularx}{\linewidth}{>{\bfseries}p{3.4cm} l X}
\toprule
\rowcolor{tabhead}\color{black}\textbf{Herramienta} &
  \color{black}\textbf{Licencia} &
  \color{black}\textbf{Función en el \textit{pipeline}} \\
\midrule
MarkItDown & MIT & Convierte PDFs a Markdown preservando estructura \\
\addlinespace
Graphify & MIT & Extrae nodos/aristas y genera grafo + \textit{wiki} \\
\addlinespace
Claude Code (CLI) & Suscripción & Entorno de consulta al grafo \\
\addlinespace
R $\geq$ 4.3 + ragnar & MIT & Normalización, \textit{chunking} e indexación RAG \\
\addlinespace
DuckDB & MIT & \textit{Vector store} local para el índice RAG \\
\addlinespace
Ollama & MIT & \textit{Embeddings} locales con bge-m3 \\
\addlinespace
Python $\geq$ 3.10 & PSF & \textit{Scripts} de descarga y conversión (\texttt{code/}) \\
\bottomrule
\end{tabularx}
\end{table}

% \begin{infobox}[Equipamiento utilizado en este artículo]
% El \textit{pipeline} completo se ejecutó en un MacBook Air M5 con 32~GB de memoria
% unificada y 1~TB SSD. No se requirió GPU dedicada: el chip Apple Silicon M5 ejecuta
% Ollama y el modelo bge-m3 de forma nativa y eficiente. La indexación del
% \textit{corpus} completo de 523~documentos tomó aproximadamente 90~minutos. El
% \textit{vector store} resultante ocupa 1{,}1~GB en disco.
% \end{infobox}

\subsection{Organización del artículo}

La Sección~\ref{sec:corpus} describe cómo construir el \textit{corpus}, desde la recolección hasta la 
conversión y normalización. La Sección~\ref{sec:grafo} documenta el procedimiento
de Graphify. La Sección~\ref{sec:explotar} explica cómo usar el grafo para
investigación. La Sección~\ref{sec:rag} presenta el sistema RAG y su funcionamienta. 
La Sección~\ref{sec:ejemplos} reproduce ejemplos reales de
consultas a ambos instrumentos. Las Secciones~\ref{sec:discusion} y
\ref{sec:conclusion} cierran con discusión y conclusiones.

% =============================================================================
\section{El \textit{corpus}: de PDFs a texto estructurado}
\label{sec:corpus}
% =============================================================================

Un \textit{corpus} bien preparado es fundamental. El mejor modelo de
\textit{embeddings} no recupera bien un documento cuyo texto está fragmentado por saltos
de página o contaminado por el ruido de la conversión de un PDF escaneado. Esta
sección describe cómo pasar de los PDFs de origen a un conjunto de archivos de
texto limpio, estructurado y listo para ser indexado.

\subsection{Composición del \textit{corpus}}

El \textit{corpus} opera en tres tracks con lógicas distintas. La Tabla~\ref{tab:track-a}
describe los Tracks~A y C, ya procesados e incorporados a ambos instrumentos.
La Tabla~\ref{tab:track-b} cubre el Track~B, pendiente de incorporación a la fecha.

\begin{table}[htbp]
\centering
\caption{Tracks A y C — \textit{corpus} procesado (520 documentos)}
\label{tab:track-a}
\footnotesize
\begin{tabular}{@{}l l r@{}}
\toprule
\rowcolor{tabhead}\color{black}\textbf{Subcarpeta en \texttt{raw-md/}} &
  \color{black}\textbf{Tipo} &
  \color{black}\textbf{N°} \\
\midrule
\multicolumn{3}{l}{\itshape Track A — Jurisprudencia} \\
\texttt{sentencias-tdlc-md/} & Sentencias TDLC & 212 \\
\texttt{sentencias-cs-md/} & Sentencias Corte Suprema & 128 \\
\texttt{resoluciones-tdlc-md/} & Resoluciones TDLC & 89 \\
\texttt{resoluciones-cs-md/} & Resoluciones Corte Suprema & 13 \\
\texttt{ae-tdlc-md/} & Acuerdos Extrajudiciales (FNE–TDLC) & 38 \\
\texttt{ern-tdlc-md/} & Expedientes de Recomendación Normativa & 22 \\
\texttt{icg-tdlc-md/} & Instrucciones de Carácter General & 9 \\
\midrule
\multicolumn{3}{l}{\itshape Track C — Normativa de base} \\
\texttt{normativa-fusiones-fne-md/} & Guías, reglamentos y formularios FNE & 8 \\
\texttt{normativa-general-md/} & DL~211 refundido & 1 \\
\midrule
\multicolumn{2}{r}{\textbf{Total procesado}} & \textbf{520} \\
\bottomrule
\end{tabular}
\end{table}

\begin{table}[htbp]
\centering
\caption{Track B — Enforcement FNE (pendiente de incorporación)}
\label{tab:track-b}
\footnotesize
\begin{tabular}{@{}l l r@{}}
\toprule
\rowcolor{tabhead}\color{black}\textbf{Subcarpeta en \texttt{raw/}} &
  \color{black}\textbf{Tipo} &
  \color{black}\textbf{N° aprox.} \\
\midrule
\texttt{concentraciones-o-integraciones/} & Decisiones FNE & 1.341 \\
\texttt{abusos-de-posicion-de-dominio/} & Investigaciones FNE & 561 \\
\texttt{acuerdos-colusorios/} & Investigaciones FNE & 193 \\
\texttt{actos-anticompetitivos-de-la-autoridad/} & Investigaciones FNE & 146 \\
\texttt{incumplimiento-de-sentencias-e-instrucciones/} & Investigaciones FNE & 146 \\
\texttt{restricciones-verticales/} & Investigaciones FNE & 36 \\
\texttt{competencia-desleal/} & Investigaciones FNE & 52 \\
\texttt{sin-categoria/} & Varios & 73 \\
\midrule
\multicolumn{2}{r}{\textbf{Total pendiente}} & \textbf{2.641} \\
\bottomrule
\end{tabular}
\end{table}

Cada track cumple un papel distinto. El Track~A es la base del análisis
jurisprudencial de este proyecto y constituye el precedente que el sistema debe
rastrear. El Track~B revela el estándar práctico de la FNE, identificando qué investiga, qué
metodologías aplica, y cuándo archiva; sus documentos fueron descargados desde
\texttt{fne.cl} y se encuentran en \texttt{raw/}, pero aún no han sido convertidos
ni incorporados al grafo ni al índice RAG. El Track~C funciona como ancla normativa. Todo
el material de jurisprudencia debe conectarse a los artículos del DL~211 que invocan.

\subsection{Recolección de documentos}

Los PDFs se descargan desde los repositorios públicos de cada institución:

\begin{itemize}
  \item \textbf{TDLC} (\texttt{tdlc.cl}): sentencias, resoluciones, acuerdos
    extrajudiciales (AE), instrucciones de carácter general (ICG), expedientes
    de recomendación normativa (ERN).
  \item \textbf{Corte Suprema} (\texttt{tdlc.cl}\footnote{En esta oportunidad, 
  el material se obtiene directamente de la página del TDLC, a pesar de que la información está
  disponible en la propia página de la CS.}): sentencias y resoluciones
    sobre recursos contra el TDLC.
  \item \textbf{FNE} (\texttt{fne.cl}): investigaciones, resoluciones de archivo,
    guías y reglamentos.
\end{itemize}

La descarga puede automatizarse con \textit{scripts} en Python disponibles en la
subcarpeta \texttt{code/} del repositorio. Los archivos se organizan en subcarpetas
temáticas bajo \texttt{raw/}, que \texttt{raw-md/} replicará tras la conversión.

\subsection{Conversión con MarkItDown}

Los PDFs no pueden procesarse directamente. Un asistente de inteligencia artificial (IA) leería ruido de
formato en lugar de contenido. MarkItDown, de Microsoft, convierte cada documento a Markdown
preservando su estructura semántica (títulos, listas, tablas) y descartando el
formato visual.\autocite{microsoft2024markitdown}

\begin{codebox}
pip3 install 'markitdown[pdf]'
python3 code/batch_convert.py       # convierte raw/ -> raw-md/
\end{codebox}

El \textit{script} detecta qué archivos ya fueron convertidos y solo procesa los nuevos
o modificados. Para PDFs escaneados, existe una variante con reconocimiento óptico de
caracteres (OCR)\footnote{Las sentencias del TDLC
anteriores a 2010 suelen caer en esta categoría; el rendimiento es desigual y
conviene revisar \texttt{ocr\_errors.log} tras la conversión.}.

\subsection{Normalización para el RAG}
\label{sub:normalizacion}

Los archivos Markdown pueden alimentar directamente a Graphify, que tolera
variaciones de formato. El RAG, en cambio, requiere texto limpio. La calidad del
\textit{chunking}, que es la división del documento en fragmentos, depende de que los
marcadores estructurales sean legibles. El \textit{corpus} jurídico chileno presenta al menos
tres problemas sistemáticos que deben resolverse antes de indexar.

Primero, la conversión de PDFs puede producir pares de espacios dentro de las líneas.
El \textit{chunker} los interpretaría como delimitadores y generaría fragmentos
artificialmente cortos.

Segundo, el OCR sobre documentos escaneados
divide los párrafos en líneas de 60–80 caracteres. Sin corrección, cada línea
podría convertirse en un fragmento independiente, perdiendo el contexto del
considerando al que pertenece.

Por último, el \textit{corpus} abarca
veintidós años (2004–2026) y refleja las convenciones tipográficas de cada período.
Los numerales que encabezan los considerandos aparecen de formas distintas: 
todo en mayúsculas (\texttt{VIGÉSIMO SEGUNDO:}), con
mayúscula inicial (\texttt{Vigésimo segundo:}) y en forma decimal. Para que el
RAG pueda citar cada considerando de forma unificada, se necesita un patrón de
normalización que cubra todos los casos.

El \textit{script} \texttt{rag/R/00\_normalize.R} resuelve los tres problemas mencionados en
un solo paso. Así, colapsa espacios dobles a uno; une las líneas de párrafo
partidas, respetando saltos de párrafo genuinos y encabezados; e inserta
prefijos \texttt{\#\#} y \texttt{\#\#\#} en los marcadores estructurales para que
el \textit{chunker} pueda asociar cada fragmento a su sección.

% =============================================================================
\section{El grafo de conocimiento}
\label{sec:grafo}
% =============================================================================

Graphify es una herramienta de código abierto que extrae conceptos y relaciones
de un \textit{corpus} de documentos y los organiza en un \emph{grafo de
conocimiento}, esto es, una red donde los nodos son conceptos jurídicos, fallos y
documentos, y las aristas representan las relaciones entre ellos.\autocite{shamsi2025graphify}
A diferencia de una búsqueda convencional, el grafo no solo encuentra documentos, sino que también
revela qué tan conectado está un concepto con el resto del campo, dónde hay vacíos
jurisprudenciales y dónde hay tensiones entre fuentes de distinta autoridad.

\subsection{El esquema de extracción: CLAUDE.md}
\label{sec:claude}

\texttt{CLAUDE.md} es el archivo de instrucciones que el investigador redacta
antes de construir el grafo y el único punto de intervención directa en la
extracción. Todo lo que contiene se incorpora al contexto de trabajo del asistente.
Su función es configurar \emph{qué} extraer de los documentos y fijar
\emph{cómo} responder a las consultas posteriores. La versión
anonimizada está en el Anexo~\ref{anx:claude-md}.

Las instrucciones específicas pueden agruparse en cuatro categorías.

\textbf{Instrucción concepto-céntrica.} Una instrucción explícita orienta la
extracción hacia los conceptos jurídicos y económicos del campo. El resultado
práctico es que los nodos más conectados pasan a ser conceptos jurídicos como
"Abuso de Posición Dominante" o "Mercado Relevante", en lugar de entidades como
"FNE" o "TDLC".

\textbf{Jerarquía de autoridad.} El \textit{corpus} contiene documentos con peso
jurídico desigual. Codificar esa jerarquía explícitamente permite que el sistema
indique si un criterio del TDLC fue confirmado, modificado o revocado por la CS,
en lugar de tratar ambas como fuentes equivalentes.

\textbf{Anclas conceptuales.} Designar los términos centrales del campo como
nodos de alta prioridad concentra la extracción y reduce el ruido. Para libre
competencia chilena, las anclas puede incluir colusión y \textit{bid rigging}, abuso
de posición dominante, precio predatorio, mercado relevante, barreras de entrada
y estándar de prueba, entre otros conceptos.

\textbf{Extracción a nivel de proposición.} Tratar cada sentencia como un único
nodo oculta su contenido. El esquema instruye al asistente a extraer las
proposiciones afirmativas de cada documento como nodos derivados. Dos sentencias
que parecen alineadas pueden defender proposiciones incompatibles; ese detalle
solo es visible a este nivel de granularidad.

\subsection{Construcción del grafo}

Graphify se ejecuta dentro de Claude Code apuntando a la carpeta de
Markdown normalizado:

\begin{codebox}
# Dentro de Claude Code, apuntando al corpus:
/graphify raw-md/
# Cuando se consulta el alcance: all
# En sesiones posteriores (retoma desde caché SHA-256):
/graphify raw-md/ --update
\end{codebox}

Graphify opera en diferentes etapas, finalizando con un HTML interactivo,
JSON consultable y un informe de auditoría en Markdown.

Para un \textit{corpus} de 520~documentos como el Track~A+C, la extracción completa
toma entre 15 y 30~minutos repartidos en una o dos sesiones. En sesiones
posteriores, solo los archivos nuevos o modificados se re-procesan.

% \begin{infobox}[Dos lecciones de operación]
% \textbf{Concurrencia.} \texttt{max\_concurrency=4} es un valor conservador que
% evita errores de límite de tasa de la API. Con cupos más altos puede subirse a
% 8 o 16.

% \textbf{Deduplicación.} Graphify puede crear nodos con etiquetas próximas para el
% mismo concepto ("mercado relevante" frente a "definición de mercado relevante").
% Una pasada de consolidación tras cada ejecución, ya sea manual o mediante un
% \textit{prompt} de limpieza a Claude Code, mejora la calidad del grafo.
% \end{infobox}

\subsection{La \textit{wiki} temática}

Una vez construido el grafo, se genera una \textit{wiki} donde cada comunidad
temática se convierte en un artículo:

\begin{codebox}
/graphify raw-md/ --wiki
\end{codebox}

Cada artículo sigue una estructura fija delimitada en \texttt{CLAUDE.md}. 
Los artículos se depositan en \texttt{graphify-out/wiki/} y pueden
consultarse en Obsidian, que renderiza los hipervínculos entre páginas y ofrece una
vista de grafo del \textit{corpus} \textit{wiki}.

Para incorporar documentos nuevos basta depositar el PDF en la subcarpeta
correspondiente y ejectuar la conversión con MarkItDown y
\texttt{/graphify raw-md/ --update}. El grafo existente se conserva.

\subsection{Métricas y hallazgos estructurales}

\begin{table}[htbp]
\centering
\caption{Métricas del grafo — Tracks A + C}
\label{tab:metricas-grafo}
\small
\begin{tabular}{ll}
\toprule
\rowcolor{tabhead}\color{black}\textbf{Métrica} & \color{black}\textbf{Valor} \\
\midrule
Documentos fuente & 520 \\
Nodos & 616 \\
Aristas & 858 \\
Comunidades & 173 \\
Artículos de \textit{wiki} generados & 184 \\
Densidad & 0{,}0045 \\
Grado medio & 2{,}79 \\
Coeficiente de clustering ($C$) & 0{,}131 \\
Longitud media de camino ($L$) & 3{,}343 \\
Correlación log-log grado/frecuencia & $-0{,}887$ \\
Coeficiente de mundo pequeño ($\sigma$) & 21{,}69 \\
Modularidad ($Q$) & 0{,}519 \\
\bottomrule
\end{tabular}
\end{table}

\noindent
La correlación log-log entre grado y frecuencia de $-0{,}887$ confirma que la distribución de grado
sigue una ley de potencia: el grafo es de tipo \emph{scale-free}, es decir,
unos pocos nodos concentran la mayoría de las conexiones. El coeficiente de
mundo pequeño $\sigma = 21{,}69$ confirma a la vez alta modularidad temática y
conexión eficiente entre comunidades. La entropía normalizada entre comunidades
(0{,}749) indica que la atención del \textit{corpus} está distribuida en numerosos
\textit{clusters}, sin un núcleo temático único dominante.

\begin{table}[htbp]
\centering
\caption{Las cinco comunidades de mayor tamaño}
\label{tab:comunidades}
\small
\begin{tabular}{clll}
\toprule
\rowcolor{tabhead}\color{black}\textbf{Rank} &
  \color{black}\textbf{Nodos} &
  \color{black}\textbf{\% del total} &
  \color{black}\textbf{Tema central} \\
\midrule
1 & 61 & 9{,}9\,\% & Operación de Concentración \\
2 & 50 & 8{,}1\,\% & Abuso de Posición Dominante \\
3 & 43 & 7{,}0\,\% & Barreras a la Entrada \\
4 & 42 & 6{,}8\,\% & Poder de Mercado \\
5 & 42 & 6{,}8\,\% & Sentencia TDLC N°~80/2009 \\
\bottomrule
\end{tabular}
\end{table}

\begin{table}[htbp]
\centering
\caption{Nodos más conectados del grafo chileno (\emph{god nodes})}
\label{tab:god-nodes}
\small
\begin{tabularx}{\linewidth}{>{\raggedright}p{5.8cm} l >{\raggedright\arraybackslash}X}
\toprule
\rowcolor{tabhead}\color{black}\textbf{Nodo (conexiones)} &
  \color{black}\textbf{Tipo} &
  \color{black}\textbf{Materia} \\
\midrule
Abuso de Posición Dominante (98) & Concepto & Conducta unilateral \\
Mercado Relevante (87) & Concepto & Definición de mercado \\
Barreras a la Entrada (72) & Concepto & Condiciones de entrada \\
Sentencia TDLC N°~47/2006 (62) & Sentencia & Fallo de referencia \\
DL~211 (53) & Normativa & Ley marco de competencia \\
Sentencia TDLC N°~80/2009 (52) & Sentencia & Fallo de referencia \\
Colusión (50) & Concepto & Acuerdos horizontales \\
\bottomrule
\end{tabularx}
\end{table}

\noindent
Los tres primeros lugares los ocupan conceptos jurídicos, lo que confirma 
el efecto de las instrucciones concepto-céntricas del
esquema.

% =============================================================================
\section{Cómo explotar el grafo para investigación}
\label{sec:explotar}
% =============================================================================

Una vez construido el grafo, el investigador puede interrogarlo en lenguaje
natural desde Claude Code con \texttt{/graphify query}. El valor de esta etapa
no está en las respuestas individuales, sino en la capacidad de hacer preguntas
que la lectura directa del \textit{corpus} no podría responder, relacionadas con la
forma del campo y no tanto sobre el contenido de un documento específico. Los ejemplos
completos figuran en la Sección~\ref{sec:ejemplos}.

\subsection{Caminos ausentes y silencios jurisprudenciales}

El grafo no solo muestra lo que el \textit{corpus} dice. También muestra lo que
\emph{no} dice. Un nodo aislado es un
documento que el sistema no pudo vincular con otros. Puede indicar que trata un
tema marginal o que la jurisprudencia sobre ese tema simplemente no ha madurado. Una
conexión ausente entre dos conceptos que deberían estar vinculados es igualmente
informativa.

Un ejemplo concreto del \textit{corpus} chileno es que el nodo "Eficiencias" existe con
grado~14, pero su conexión con sentencias donde el argumento fue acogido para
mitigar una sanción es débil. El grafo señala una brecha; no es que el
\textit{corpus} ignore las eficiencias, sino que no las conecta con resultados
favorables (o desfavorables) de forma sistemática. Eso puede ser el punto de partida de una agenda
de investigación o de una estrategia de distinción en litigio.

\subsection{Tensiones de autoridad TDLC/CS}

El grafo puede rastrear de forma sistemática los casos donde la CS confirmó o
revocó, total o parcialmente, sentencias o resoluciones del TDLC. Esto no es
posible mediante búsqueda convencional salvo que el investigador ya sepa qué
casos buscar.

% =============================================================================
\section{El sistema RAG vectorial}
\label{sec:rag}
% =============================================================================

El RAG (\emph{Retrieval-Augmented Generation}) permite formular búsquedas semánticas y 
encontrar el texto exacto dentro del \textit{corpus} en que se define un término
o se establece un estándar.

\subsection{\textquestiondown Qué hace un RAG?}

Un RAG tiene dos fases. La primera, \textit{offline}, divide el \textit{corpus} en fragmentos
pequeños, los convierte en representaciones numéricas (\emph{embeddings}) que
capturan su significado, y los almacena en una base de datos de búsqueda. La
segunda, en consulta, toma la pregunta del usuario, la convierte a la misma
representación numérica, recupera los fragmentos más similares del índice, y se
los entrega a un modelo de lenguaje para que genere una respuesta citando la
fuente exacta.

El modelo de lenguaje solo ve los fragmentos recuperados, no el \textit{corpus}
completo. Esto controla el costo, reduce las fabricaciones de información y hace
auditable la fuente de cada afirmación.

\subsection{El entorno: ragnar + DuckDB}

El sistema usa el paquete \texttt{ragnar} de R, con DuckDB como
\textit{vector store} local.\footnote{La implementación está basada en el tutorial de
Bastián Olea (\url{https://bastianolea.rbind.io/blog/rag_ragnar/}), que documenta
el patrón correcto de uso de ragnar~v2.} En esta etapa, \texttt{ragnar} divide el texto en
fragmentos según la estructura del Markdown, construye los índices de búsqueda y
los almacena en DuckDB. El índice es un archivo único portable.

El mismo entorno R que ejecuta la normalización también ejecuta el RAG, sin
cambiar de lenguaje ni de infraestructura.

% \begin{infobox}[Nota técnica: ragnar v2 y el atributo \texttt{@document@origin}]
% La versión~2 del \textit{store} de ragnar almacena el origen de cada fragmento en
% el atributo S7 \texttt{@document@origin} del objeto
% \texttt{MarkdownDocumentChunks}, no en columnas del tibble resultante. Es
% necesario crear el documento con \texttt{ragnar::read\_as\_markdown(path)} —que
% fija el origen automáticamente— en lugar de leer el texto directamente. Sin esto,
% el campo de origen de todos los fragmentos resulta vacío, inutilizando la
% trazabilidad de las citas. Este detalle no está documentado en ragnar; se
% identificó leyendo el código fuente de \texttt{ragnar\_store\_insert\_v2}.
% \end{infobox}

\subsection{El modelo de representación semántica: bge-m3}

Un modelo de \textit{embedding} convierte texto en un vector numérico de alta
dimensión. La calidad de esa conversión determina la calidad de la recuperación.
Se eligió bge-m3 (BAAI) por tres razones relevantes para este \textit{corpus}.

Primero, opera en más de cien idiomas, incluyendo el español
jurídico técnico del \textit{corpus} chileno. Segundo, procesa grandes bloques de palabras 
de una sola vez, lo que permite mantener la coherencia de un
considerando largo sin necesidad de partirlo. Tercero, corre completamente en el
equipo del investigador vía Ollama, sin enviar datos a servicios externos y
sin costo de API.

\begin{codebox}
ollama serve                 # en una pestaña separada del terminal
ollama pull bge-m3           # solo la primera vez (~1.2 GB)
\end{codebox}

\subsection{Construcción del índice}

\begin{codebox}
Rscript rag/R/00_normalize.R    # ya ejecutado en la seccion anterior
Rscript rag/R/01_chunk_index.R  # chunking + bge-m3 + store.duckdb
\end{codebox}

El paso siguiente es aplicar una función que divide cada documento normalizado
en fragmentos usando los encabezados Markdown como señales estructurales. Cada
fragmento lleva el campo \texttt{context} con la jerarquía de secciones que lo
precede, lo que permite al modelo citar con precisión el número del considerando.

Se construyen dos índices en paralelo:

\begin{itemize}
  \item \textbf{FTS (BM25):} índice léxico de texto completo. Siempre disponible,
    sin dependencias externas. Robusto para nombres de partes, números de rol y
    artículos del DL~211.
  \item \textbf{VSS (vectorial):} índice de similitud semántica sobre los
    \textit{embeddings} de bge-m3.
\end{itemize}

\subsection{Sistema de instrucciones}

El archivo \texttt{rag/R/prompt\_rules.R} inyecta como \textit{system prompt} las
mismas reglas de dominio codificadas en \texttt{CLAUDE.md}: jerarquía de fuentes,
obligación de citar considerandos textuales, y
la regla de tensión TDLC/CS (toda cita del TDLC debe indicar si la CS confirmó,
modificó, revocó o no conoció el punto). Esto garantiza que las mismas reglas que
gobiernan el grafo también gobiernen el RAG.

% \subsection{Interfaz de consulta}

% \begin{codebox}
% ./rag/ask.sh "¿Qué estándar aplicó el TDLC para condenar por colusión?"
% ./rag/ask.sh "tu pregunta" 20   # top_k=20 para documentos largos
% \end{codebox}

% El parámetro \texttt{top\_k} controla cuántos fragmentos del \textit{corpus} se
% entregan al modelo para generar la respuesta (por omisión: 12). Un valor más alto
% aumenta la probabilidad de capturar fragmentos al final de documentos largos
% —como el dispositivo resolutivo de una sentencia CS ubicado en la última página—
% a costo de un procesamiento ligeramente mayor. El modelo de generación puede
% cambiarse en \texttt{rag/R/config.R}:

% \begin{table}[htbp]
% \centering
% \caption{Modelos de generación por tipo de tarea}
% \label{tab:gen-models}
% \small
% \begin{tabularx}{\linewidth}{>{\bfseries}p{4cm} l X}
% \toprule
% \rowcolor{tabhead}\color{black}\textbf{Modelo} &
%   \color{black}\textbf{Variable} &
%   \color{black}\textbf{Uso recomendado} \\
% \midrule
% claude-haiku-4-5 & \texttt{GEN\_FAST} &
%   Extracción de considerandos, Q\&A simple \\
% \addlinespace
% claude-sonnet-4-6 & \texttt{GEN\_SMART} &
%   Minutas, análisis TDLC/CS, contradicciones \\
% \addlinespace
% claude-opus-4-8 & \texttt{GEN\_TOP} &
%   Síntesis para paper, registro de hipótesis \\
% \bottomrule
% \end{tabularx}
% \end{table}

% \noindent
% Para uso interactivo, \texttt{ask.sh} usa \texttt{claude --print} autenticado
% mediante suscripción a Claude.ai, sin clave de API separada.

% =============================================================================
\section{Ejemplos de consulta}
\label{sec:ejemplos}
% =============================================================================

Los ejemplos siguientes reproducen consultas reales ejecutadas sobre el
\textit{pipeline} documentado. \textbf{Las respuestas fueron abreviadas}.

\subsection{Graphify: consultas estructurales}

\subsubsection*{Ejemplo 1 --- Conceptos con más conexiones}

\begin{codebox}
/graphify query "qué nodos de concepto tienen más conexiones?"
\end{codebox}

\begin{quote}
\itshape
Los conceptos con mayor grado son Abuso de Posición Dominante (98~conexiones),
Mercado Relevante (87), Barreras a la Entrada (72), DL~211 (53) y Colusión
(50).[\ldots]
\end{quote}

\subsubsection*{Ejemplo 2 --- Fallos con mayor alcance transversal}

\begin{codebox}
/graphify query "qué sentencias o resoluciones conectan más de dos
comunidades temáticas distintas? cuáles son los fallos con mayor
alcance transversal en el corpus?"
\end{codebox}

\begin{quote}
\itshape
Los documentos con mayor alcance transversal son la Sentencia TDLC
N\textdegree~47/2006 (11~comunidades externas, grado~62) y la Sentencia TDLC
N\textdegree~80/2009 (11~comunidades externas, grado~52). Ambas alcanzan la
misma frontera de 11~comunidades, posicionándose como los fallos con mayor
efecto estructurador del corpus. El Acuerdo Extrajudicial N\textdegree~2/2010
queda en tercer lugar con 9~comunidades externas.

La coincidencia entre los nodos de mayor grado y los de mayor alcance
comunitario no es trivial: indica que estos fallos no solo son muy citados,
sino que lo son desde ángulos temáticos distintos. Un fallo puede ser
referencia frecuente dentro de su propia materia sin cruzar hacia otros
clusters; estos dos cruzan once. [\ldots]
\end{quote}

\subsection{RAG vectorial: recuperación con cita literal}

\subsubsection*{Ejemplo 3 --- Estrangulamiento de márgenes}

\begin{codebox}
./rag/ask.sh "qué requisitos establece el TDLC para condenar por
estrangulamiento de márgenes? Cita los considerandos exactos e indica si la
CS confirmó o modificó ese estándar."
\end{codebox}

\begin{quote}
\itshape
La Sentencia TDLC N\textdegree~156/2017 (Considerando 93\textdegree) estableció
cinco requisitos copulativos para configurar un estrangulamiento de márgenes:
(i)~presencia simultánea en el mercado mayorista y minorista; (ii)~diferencial de
precios que haga imposible o antieconómica la operación de un competidor igualmente
eficiente; (iii)~nexo causal entre la conducta y el efecto exclusorio; (iv)~poder
de mercado en el segmento mayorista; y (v)~ausencia de justificación objetiva de la
estructura de precios. En el Considerando~97\textdegree~se aplicó el test CIE
(Competidor Igualmente Eficiente) como metodología de evaluación.

La Sentencia N\textdegree~209/2026 reprodujo y aplicó el mismo estándar de cinco
requisitos, confirmando su continuidad hacia la fecha de este artículo.

[Respecto de la CS:] La sentencia N\textdegree~156/2017 no fue recurrida en este
punto; el TDLC conoció el caso en revisión de fallo arbitral. No consta
pronunciamiento de la CS sobre los cinco requisitos como tal. [\ldots]
\end{quote}

\subsubsection*{Ejemplo 4 --- CDF: multa y cálculo de la sanción}

\begin{codebox}
./rag/ask.sh "de cuánto fue la multa impuesta a CDF? Cómo se calculó?"
\end{codebox}

\begin{quote}
\itshape
La Sentencia TDLC N\textdegree~191/2024 impuso en su Resolutivo N\textdegree~5 una
multa de 32.000 Unidades Tributarias Anuales (UTA) a CDF por infracción al
artículo~3\textdegree~del DL~211. Para la graduación, el Tribunal aplicó los
criterios del artículo~26 del DL~211, ponderando: (i)~el beneficio económico
obtenido por la conducta anticompetitiva, estimado a partir de los ingresos
generados por las restricciones de competencia en el mercado de transmisión de
señales deportivas de pago; (ii)~la gravedad de la infracción y su extensión
temporal; y (iii)~la afectación al interés general. El Tribunal no apreció
circunstancias atenuantes que justificaran una reducción y estimó el monto
proporcional a la gravedad de la conducta. [\ldots]

[Regla de tensión TDLC/CS:] La Corte Suprema, al conocer el recurso, confirmó
la sanción en el Considerando~19\textdegree: \emph{"el monto de la multa
impuesta a CDF aparece como adecuado y proporcional a la gravedad de la conducta
y al beneficio económico obtenido, por lo que no se advierte mérito para
modificarla"}. No existe tensión entre instancias en este punto. [\ldots]
\end{quote}

% =============================================================================
\section{Discusión}
\label{sec:discusion}
% =============================================================================

\subsection{Investigación estructural y práctica diaria}

La separación entre grafo y RAG responde a necesidades del usuario
genuinamente distintas. El abogado que prepara un escrito necesita el texto exacto
del considerando que respalda un argumento; el RAG resuelve esa tarea. Cuando el
mismo profesional diseña la estrategia de un caso nuevo y necesita saber, por
ejemplo, cuántas veces se rechazó un argumento específico, el grafo es
la herramienta adecuada.

Un sistema único que pretenda satisfacer ambas necesidades —típicamente, un
\textit{chatbot} con búsqueda semántica— termina sirviendo parcialmente a los dos
sin optimizar a ninguno. La arquitectura propuesta aquí separa instrumentos y les
asigna las preguntas para las que están diseñados.

\subsection{El valor de las brechas}

La búsqueda convencional entrega lo que el \textit{corpus} dice sobre un tema. Lo que
no entrega es el mapa de lo que el \textit{corpus} \emph{no} dice. El grafo lo produce
como subproducto natural. Las brechas de conexión, los nodos aislados y los
conceptos invocados pero no definidos revelan las zonas del campo donde no existe
criterio jurisprudencial consolidado. Para el investigador de libre competencia, esa información
alimenta tanto la agenda académica como la estrategia de litigio.

\subsection{Control del error}

El riesgo de fabricar información existe en todo sistema de IA, pero es
considerablemente menor en este diseño que en un \textit{chatbot} de propósito
general. El grafo restringe al asistente al contenido extraído del \textit{corpus}.
El RAG solo puede afirmar lo que está en los fragmentos recuperados; si el
fragmento no está, el modelo lo declara explícitamente. Con todo, siempre conviene
verificar los hallazgos contra los documentos fuente, especialmente para citas
textuales largas.

\subsection{Extensiones}

La incorporación del Track~B (2.641 documentos FNE) es la extensión más directa:
expande el grafo y amplía el \textit{retrieval} a las resoluciones de archivo y las
condiciones impuestas en concentraciones.

% =============================================================================
\section{Conclusión}
\label{sec:conclusion}
% =============================================================================

Este artículo documentó la construcción de dos instrumentos complementarios para
la investigación y la práctica en libre competencia chilena: un grafo de
conocimiento concepto-céntrico sobre 520~documentos y un sistema RAG vectorial
sobre el mismo \textit{corpus}. Ambos son replicables, de código abierto y
ajustables a cualquier \textit{corpus} jurídico.

Las contribuciones son específicas. Primero, el diseño del
esquema en \texttt{CLAUDE.md} —jerarquía CS/TDLC/FNE, instrucción
concepto-céntrica, anclas conceptuales, extracción a nivel de proposición— es
lo que transforma la materia prima en un mapa jurisprudencial. Sin ese diseño previo,
el grafo produce lo que resulta más prominente en los documentos, que puede no ser 
útil para la investigación. Segundo, el diseño dual 
responde a necesidades genuinamente distintas del profesional que litiga e investiga.

El repositorio está disponible públicamente con el \texttt{CLAUDE.md} anotado,
los \textit{scripts} de descarga y conversión, el código RAG completo y este
artículo, para que investigadores de otras jurisdicciones o áreas del derecho
adapten el \textit{pipeline} a su propio \textit{corpus}.

La invitación es a usar el sistema, mejorar el esquema, identificar
los \emph{god nodes} del propio campo e interrogar sus silencios.

% =============================================================================
% REFERENCIAS
% =============================================================================

\printbibliography

% =============================================================================
% ANEXOS
% =============================================================================

\appendix

\section{Contenido del archivo CLAUDE.md (versión anonimizada)}
\label{anx:claude-md}

El archivo \texttt{CLAUDE.md} es el esquema de diseño que moldea la extracción.
A continuación se reproduce su estructura con las rutas absolutas y la información
personal removidas. Los elementos en \texttt{[...]} son \emph{placeholders}
para que el lector complete con sus propios datos. La Sección~\ref{sec:claude}
discute las cuatro decisiones de diseño que distinguen este esquema.

\begin{codebox}
# Base de Conocimiento — [NOMBRE DEL CAMPO]

## Qué es esto
Una base de conocimiento de [práctica profesional / investigación] sobre
[descripción del dominio]. El corpus cubre [tipos de documentos].
El sistema está orientado a [descripción del usuario].

## Cómo está organizado
- raw/         — documentos fuente en PDF. No modificar.
- raw-md/      — versiones Markdown de raw/. No modificar.
- wiki/        — wiki organizada por temas. La IA la mantiene.
- notes/       — notas personales del investigador. No son fuentes
                 de autoridad y nunca deben citarse como tal.
- outputs/     — minutas, informes y documentos generados.
- code/        — scripts de descarga y transformación.

## Jerarquía de autoridad de las fuentes
1. [Fuente de mayor autoridad] — vinculante
2. [Segunda fuente] — especializado; supeditado a 1
   2b. [Variante normativa de 2] — vinculante para agentes del mercado
   2c. [Acuerdos entre instituciones] — vinculantes para las partes
3. [Autoridad administrativa]
4. [Guías no vinculantes]
5. [Estudios sectoriales]
6. [Informes periciales]
7. [Doctrina académica]
0. Notas personales — nivel cero; sin autoridad jurídica

## Regla de conflicto
Cuando fuentes de diferentes niveles aborden el mismo punto, citar
la de mayor autoridad y señalar cómo confirma, modifica o anula
a la de menor autoridad.

## Ponderación por centralidad de nodo (Node Centrality Weighting)
El grafo rankea los documentos por conexiones. Los más conectados
son "god nodes" y reflejan el consenso del campo. Al generar artículos
de wiki y sintetizar respuestas, dar mayor peso a los god nodes.

## Reglas de la wiki
- Cada tema tiene su propio .md en wiki/
- Estructura fija: (1) Resumen, (2) Jurisprudencia, (3) Enforcement,
  (4) Casos vinculados, (5) Lo que las fuentes no abordan
- Vincular temas con [[nombre-del-tema]]
- Mantener INDEX.md con descripción de una línea por tema
- Al agregar nuevas fuentes, actualizar los artículos relevantes

## Conceptos centrales (anclas de extracción — extraer como nodos
## de alta prioridad)
[Lista de 15-30 conceptos centrales del campo]

## Reglas de extracción de proposiciones
Para cada documento, extraer proposiciones afirmativas con:
- Premisa: hecho empírico con referencia al párrafo
- Inferencia: conclusión jurídica/económica
- Puntero: referencia exacta al texto

## Registro de hipótesis
Cada vez que una consulta revele una contradicción, generar una
hipótesis y registrarla en hypotheses.md con: fecha, consulta de
origen, documentos involucrados, tipo de inconsistencia.

## Módulo de diagnóstico
Después de cada build, calcular propiedades de red del grafo y
reportarlas en GRAPH_DIAGNOSTICS_METRICS.md:
- Distribución de grado (¿sigue una ley de potencia?)
- Coeficiente de clustering + longitud media de camino
- Entropía entre comunidades
- Traza año a año

## Perspectiva de uso
[Descripción del rol del usuario y sus necesidades principales]

## Reglas de respuesta
1. Citas textuales — obligatorio (texto exacto + referencia precisa)
2. Nivel de detalle — intermedio aplicado
3. Sin precedente local — indicarlo explícitamente
4. Ventana temporal — sin corte fijo; último pronunciamiento relevante
5. Precedentes recientes — énfasis en los últimos 5-7 años
\end{codebox}

% =============================================================================
\section{Glosario de términos técnicos}
\label{anx:glosario}
% =============================================================================

\begin{description}[leftmargin=3.2cm, style=sameline, itemsep=6pt]

\item[bge-m3] Modelo de representación semántica desarrollado por el Beijing
Academy of Artificial Intelligence (BAAI). Es multilingüe (más de cien idiomas),
soporta bloques de texto de hasta 8.192 palabras y puede ejecutarse localmente
vía Ollama sin costo de API.

\item[BM25] Algoritmo de búsqueda léxica basado en la frecuencia de las palabras
en los documentos. Permite encontrar fragmentos relevantes sin representaciones
matemáticas sofisticadas. Es robusto para términos técnicos específicos como
nombres de partes o números de artículos.

\item[\textit{chunking}] División de un documento largo en fragmentos más pequeños para
que el modelo de lenguaje pueda procesarlos. Cada fragmento incluye información
sobre su contexto (por ejemplo, a qué sección del fallo pertenece).

\item[DuckDB] Base de datos que se instala en el computador del investigador,
sin requerir un servidor externo. Almacena los fragmentos del \textit{corpus} y los
índices de búsqueda en un solo archivo portable.

\item[\textit{embedding} / vector semántico] Representación numérica de un texto que
captura su significado. Dos textos con significados similares tendrán
representaciones cercanas, lo que permite encontrar documentos relacionados
aunque usen palabras distintas.

\item[Graphify] Herramienta de código abierto que usa inteligencia artificial
para extraer conceptos y relaciones de un \textit{corpus} de documentos y construir un
grafo de conocimiento. Desarrollada por Bita Shamsi y colaboradores.

\item[grafo de conocimiento] Red de nodos conectados por aristas. En este
trabajo, los nodos son conceptos jurídicos, fallos y documentos; las aristas
representan relaciones entre ellos. El grafo permite navegar el \textit{corpus} de forma
estructural.

\item[god node] Nodo con un número muy alto de conexiones en el grafo. En un
\textit{corpus} jurídico, los god nodes son los conceptos o fallos que el campo trata
como referencia fundamental y que emergen de los patrones de citación.

\item[MarkItDown] Herramienta de Microsoft que convierte archivos PDF a texto
en formato Markdown, preservando la estructura del documento (títulos, listas,
tablas) y descartando el formato visual.

\item[Markdown] Formato de texto simple que usa caracteres especiales (\#, *,
etc.) para indicar estructura como títulos, listas y énfasis. Es legible por
humanos y por sistemas de inteligencia artificial.

\item[Ollama] Aplicación que permite ejecutar modelos de lenguaje de forma
local, en el computador del investigador, sin enviar datos a servicios externos.

\item[ragnar] Paquete de R desarrollado por Posit que facilita la construcción
de sistemas RAG. Incluye herramientas para dividir documentos, construir índices
de búsqueda y recuperar fragmentos relevantes.

\item[RAG (\textit{Retrieval-Augmented Generation})] Técnica que combina búsqueda y
generación de texto. Ante una pregunta, el sistema primero recupera los
fragmentos del \textit{corpus} más relevantes y luego se los entrega a un modelo de
lenguaje para que genere una respuesta con citas textuales.

\item[\textit{scale-free}] Propiedad de una red en la que unos pocos nodos concentran
la mayoría de las conexiones (los god nodes) mientras la mayoría tiene muy pocas.
Es la firma estructural típica de redes donde el conocimiento consolidado atrae
referencias adicionales.

\item[\textit{small-world}] Propiedad de una red en la que cualquier par de nodos está
conectado a través de pocos pasos intermedios. En el grafo chileno, cualquier
par de conceptos está separado en promedio por 3{,}3 pasos.

\item[VSS (\textit{Vector Similarity Search})] Búsqueda por similitud semántica
usando representaciones vectoriales. Complementa la búsqueda léxica (BM25) al
encontrar documentos relevantes aunque usen vocabulario distinto al de la consulta.

\end{description}

\end{document}
